\chapter{Introduzione}


Nella stesura  della tesi, in quasi tutti i casi tre parti devono emergere con chiarezza:
\begin{itemize}
	\item analisi dello stato dell’arte e definizione del problema
	\item studio e/o sviluppo di modelli per la risoluzione del problema
	\item risultati teorici o sperimentali ottenuti
\end{itemize} 



La scelta di come strutturare l'elaborato finale non è univoca, pertanto si suggerisce di iniziare creando un indice organizzato in maniera gerarchica suddiviso in capitoli, sezioni, sottosezioni (non eccedere con la profondità, si rischia di frammentare troppo il testo).

Un esempio di 'albero' che prevede annidamenti fino alle sottosezioni è il seguente (puramente indicativo):

\section{Il problema affrontato}
	\subsection{subsection a}
	\subsection{subsection b}
	
	
\section{Stato dell'arte}
\subsection{subsection c}
\subsection{subsection d}

\section{Soluzioni in letteratura}

\subsection{subsection e}
\subsection{subsection f}
\subsection{subsection g}